\chapter{Bilder, Urbilder und Fasern}
\label{b05guide:ch:images}

\section{Grundbegriffe und elementare Regeln}

\begin{definition}[Bild, Urbild und Faser]
\label{b05guide:def:images}
Für $f\colon A\to B$, $U\subseteq A$, $V\subseteq B$ und $b\in B$ sei
\[
 \begin{aligned}
 f[U]&=\{y\in B\mid\exists x\in U\,(y=f(x))\},\\
 f^{-1}[V]&=\{x\in A\mid f(x)\in V\},\\
 \operatorname{Fib}_f(b)&=f^{-1}[\{b\}].
 \end{aligned}
\]
$f[U]$ heißt \emph{Bild} von $U$, $f[A]$ die \emph{Bildmenge}
der Funktion; $f^{-1}[V]$ heißt \emph{Urbildmenge} von $V$.
$\operatorname{Fib}_f(b)$ ist die \emph{Faser über $b$}.
Die Notation $f^{-1}[V]$ bezeichnet keine Umkehrfunktion.
\end{definition}

In diesem Kapitel sei stets $f\colon A\to B$,
$U,W\subseteq A$ und $V,Z\subseteq B$, soweit ein Satz
keine anderen Voraussetzungen nennt.

\begin{theorem}[Grundmengen, leere Mengen und einzelne Werte]
\label{b05guide:thm:imagebasic}
Es gelten
\[
 \begin{gathered}
 f[U]\subseteq B,\qquad f^{-1}[V]\subseteq A,\\
 f[\varnothing]=\varnothing,\qquad
 f^{-1}[\varnothing]=\varnothing,\qquad f^{-1}[B]=A,\\
 f[\{x\}]=\{f(x)\}\quad(x\in A),\\
 f[U]=\varnothing\quad\Longleftrightarrow\quad U=\varnothing,\\
 f^{-1}[V]=\varnothing\quad\Longleftrightarrow\quad V\cap f[A]=\varnothing.
 \end{gathered}
\]
\end{theorem}

\begin{theorem}[Verträglichkeit mit Mengeninklusionen]
\[
 \begin{aligned}
 U\subseteq W&\ \Longrightarrow\ f[U]\subseteq f[W],\\
 V\subseteq Z&\ \Longrightarrow\ f^{-1}[V]\subseteq f^{-1}[Z].
 \end{aligned}
\]
\end{theorem}

\begin{theorem}[Fasern beschreiben alle Argumente und den Graphen]
\label{b05guide:thm:fibers}
Für $b,c\in B$ mit $b\ne c$ gilt
\[
 \operatorname{Fib}_f(b)\cap\operatorname{Fib}_f(c)=\varnothing.
\]
Außerdem gelten
\[
 \begin{gathered}
 A=\bigcup_{b\in B}\operatorname{Fib}_f(b),\qquad
 b\in f[A]\ \Longleftrightarrow\ \operatorname{Fib}_f(b)\ne\varnothing,\\
 f=\bigcup_{b\in B}\bigl(\operatorname{Fib}_f(b)\times\{b\}\bigr).
 \end{gathered}
\]
Leere Fasern sind erlaubt.
\end{theorem}

\begin{theorem}[Rekonstruktion aus vorgegebenen Fasern]
\label{b05guide:thm:fiberconstruction}
Sei $(E_b)_{b\in B}$ eine gegebene Familie von Teilmengen von $A$ mit
\[
 \bigcup_{b\in B}E_b=A,\qquad
 E_b\cap E_c=\varnothing\quad(b,c\in B,\ b\ne c).
\]
Dann existiert genau eine Funktion $f\colon A\to B$ mit
$\operatorname{Fib}_f(b)=E_b$ für alle $b\in B$. Ihr Graph lautet
\[
 f=\bigcup_{b\in B}(E_b\times\{b\}).
\]
\end{theorem}

\section{Bilder von Mengenoperationen}

\begin{theorem}[Bilder beliebiger Vereinigungen]
\label{b05guide:thm:imageunion}
Für eine beliebige gegebene Familie $(U_i)_{i\in I}$ von
Teilmengen von $A$ gilt
\[
 f\left[\bigcup_{i\in I}U_i\right]
   =\bigcup_{i\in I}f[U_i].
\]
Das gilt auch für $I=\varnothing$. Insbesondere ist
$f[U\cup W]=f[U]\cup f[W]$.
\end{theorem}

\begin{theorem}[Bilder von Durchschnitten und Differenzen]
\label{b05guide:thm:imageinclusions}
Es gelten im allgemeinen nur die Inklusionen
\[
 \begin{aligned}
 f[U\cap W]&\subseteq f[U]\cap f[W],\\
 f[U]\setminus f[W]&\subseteq f[U\setminus W].
 \end{aligned}
\]
Für jede nichtleere Indexmenge $I$ und $U_i\subseteq A$ gilt
\[
 f\left[\bigcap_{i\in I}U_i\right]
   \subseteq\bigcap_{i\in I}f[U_i].
\]
\end{theorem}

\section{Urbilder von Mengenoperationen}

\begin{definition}[Durchschnitt relativ zu einer Grundmenge]
Für $U_i\subseteq A$ sei
\[
 \bigcap_{i\in I}^{A}U_i
   =\{x\in A\mid\forall i\in I\,(x\in U_i)\}.
\]
Bei $I=\varnothing$ ist dieser Durchschnitt $A$.
Bei $I\ne\varnothing$ ist er der übliche Durchschnitt.
Die entsprechende Schreibweise mit Grundmenge $B$ wird ebenso verwendet.
\end{definition}

\begin{theorem}[Urbilder beliebiger Vereinigungen und Durchschnitte]
\label{b05guide:thm:preimagefamilies}
Für jede gegebene Familie $(V_i)_{i\in I}$ von Teilmengen von $B$ gilt
\[
 \begin{aligned}
 f^{-1}\left[\bigcup_{i\in I}V_i\right]
   &=\bigcup_{i\in I}f^{-1}[V_i],\\
 f^{-1}\left[\bigcap_{i\in I}^{B}V_i\right]
   &=\bigcap_{i\in I}^{A}f^{-1}[V_i].
 \end{aligned}
\]
Beide Formeln schließen die leere Indexmenge ein.
\end{theorem}

\begin{theorem}[Urbilder von Komplementen und Differenzen]
\label{b05guide:thm:preimageboolean}
Es gelten
\[
 \begin{aligned}
 f^{-1}[B\setminus V]&=A\setminus f^{-1}[V],\\
 f^{-1}[V\setminus Z]&=f^{-1}[V]\setminus f^{-1}[Z],\\
 f^{-1}[V\mathbin{\triangle}Z]
   &=f^{-1}[V]\mathbin{\triangle}f^{-1}[Z].
 \end{aligned}
\]
Dabei ist $V\mathbin{\triangle}Z=(V\setminus Z)\cup(Z\setminus V)$
die symmetrische Differenz.
\end{theorem}

\begin{theorem}[Urbilder als Vereinigungen von Fasern]
\label{b05guide:thm:preimagefibers}
Für $V\subseteq B$ gilt
\[
 f^{-1}[V]=\bigcup_{b\in V}\operatorname{Fib}_f(b).
\]
Für $f,g\colon A\to B$ gilt außerdem
\[
 f=g\quad\Longleftrightarrow\quad
 \forall b\in B\,\bigl(f^{-1}[\{b\}]=g^{-1}[\{b\}]\bigr).
\]
\end{theorem}

\section{Zusammenspiel von Bild und Urbild}

\begin{theorem}[Grundvergleich zwischen Bild und Urbild]
\label{b05guide:thm:adjunction}
\[
 f[U]\subseteq V\quad\Longleftrightarrow\quad U\subseteq f^{-1}[V].
\]
Insbesondere gilt $U\subseteq f^{-1}[f[U]]$.
\end{theorem}

\begin{theorem}[Bild einer Urbildmenge]
\label{b05guide:thm:imagepreimage}
\[
 f[f^{-1}[V]]=V\cap f[A].
\]
Die erreichten Elemente von $V$ hängen also nur von
$V\cap f[A]$ ab. Genauer gilt
\[
 f^{-1}[V]=f^{-1}[Z]
 \quad\Longleftrightarrow\quad
 V\cap f[A]=Z\cap f[A].
\]
\end{theorem}

\begin{theorem}[Mischen von Bild und Urbild]
\label{b05guide:thm:mixed}
\[
 \begin{aligned}
 f[U\cap f^{-1}[V]]&=f[U]\cap V,\\
 f[U\setminus f^{-1}[V]]&=f[U]\setminus V.
 \end{aligned}
\]
Diese Gleichheiten gelten für jede Funktion.
\end{theorem}

\begin{theorem}[Urbild der Bildmenge einer Teilmenge]
\label{b05guide:thm:preimageimage}
\[
 f^{-1}[f[U]]
   =\{x\in A\mid\exists u\in U\,(f(x)=f(u))\}.
\]
Die Menge ist die Vereinigung aller Fasern, die $U$ treffen.
Es gilt
\[
 U=f^{-1}[f[U]]
 \quad\Longleftrightarrow\quad
 \forall b\in B\,
 \bigl(\operatorname{Fib}_f(b)\cap U\ne\varnothing
       \Rightarrow\operatorname{Fib}_f(b)\subseteq U\bigr).
\]
\end{theorem}

\begin{theorem}[Wiederholtes Bild- und Urbildnehmen]
\label{b05guide:thm:repeated}
\[
 \begin{aligned}
 f[f^{-1}[f[U]]]&=f[U],\\
 f^{-1}[f[f^{-1}[V]]]&=f^{-1}[V].
 \end{aligned}
\]
\end{theorem}
